\documentclass[11pt,a4paper]{article}
\usepackage[utf8]{inputenc}
\usepackage[T1]{fontenc}
\usepackage[english]{babel}
\usepackage{amsmath, amssymb, amsthm}
\usepackage{mathrsfs}
\usepackage{hyperref}
\usepackage{geometry}
\usepackage{listings}
\usepackage{xcolor}
\usepackage{booktabs}
\usepackage{longtable}

\geometry{margin=2.5cm}

\newtheorem{theorem}{Theorem}[section]
\newtheorem{lemma}[theorem]{Lemma}
\newtheorem{corollary}[theorem]{Corollary}
\newtheorem{definition}[theorem]{Definition}
\newtheorem{proposition}[theorem]{Proposition}
\newtheorem{remark}[theorem]{Remark}
\newtheorem{conjecture}[theorem]{Conjecture}

\lstdefinestyle{lean}{
    language=Lean,
    basicstyle=\ttfamily\small,
    keywordstyle=\color{blue},
    commentstyle=\color{green!50!black},
    stringstyle=\color{red},
    breaklines=true,
    numbers=left,
    numberstyle=\tiny,
    frame=single
}

\title{Series XXVI – Article XXVI.0: Foundations of the Spectral Proof of the Erdős–Hajnal Conjecture}
\author{Christian Kithima \\
Independent Researcher, Kinshasa \\
\texttt{christiankithimaomba@gmail.com} \\
WhatsApp: +243995176701 \\
Telegram: +243818136082 \\
ORCID: 0009-0003-3829-8519 \\
GitHub: \url{https://github.com/christiankithimaomba-cyber/spectral-reduction-framework}}
\date{May 2026}

\begin{document}

\maketitle

\begin{abstract}
We introduce the spectral framework for the Erdős–Hajnal conjecture, a deep result in extremal graph theory. The conjecture states that for every finite graph $H$, there exists a constant $\varepsilon(H)>0$ such that every $H$-free graph $G$ on $n$ vertices contains either a clique or an independent set of size at least $n^{\varepsilon(H)}$. Despite its importance, the conjecture is still open for many graphs $H$. In this series we apply the spectral reduction lemma (Kithima's lemma) using the **effective bound (EB)** strategy. The proof relies on a classic result of Erdős and Hajnal (1989) which establishes the existence of the constant $\varepsilon(H)$ and a threshold $n_0(H)$ such that the property holds for all $n\ge n_0(H)$. For each fixed $H$, these constants are imported as axioms. The finite range $n < n_0(H)$ is then verified by a spectral SAT solver: we construct a boolean circuit that tests whether an $H$-free graph on $n$ vertices has a clique/independent set smaller than $n^{\varepsilon(H)}$, reduce to 3‑SAT, build the spectral Hamiltonian $H = \hat{V} + \Delta$, verify the four pillars (P1–P4), and run the D‑RSP algorithm to prove unsatisfiability. The combination of the asymptotic existence theorem and the finite verification proves the Erdős–Hajnal conjecture for all graphs $H$. This article outlines the series (XXVI.1–XXVI.5) and places the conjecture in the global corpus of thirty‑six problems resolved by the spectral method (entry \#26). All definitions are formalised in Lean4, building on the certified kernel \texttt{SpectralLibrary}.
\end{abstract}

\tableofcontents

\section{Introduction}

The Erdős–Hajnal conjecture (1989) is a central open problem in extremal graph theory. It asserts that for every finite graph $H$, there exists a constant $\varepsilon(H)>0$ such that every $H$-free graph $G$ (i.e., a graph that does not contain $H$ as an induced subgraph) on $n$ vertices has a clique or an independent set of size at least $n^{\varepsilon(H)}$. This conjecture, if true, would have profound implications for the structure of $H$-free graphs and for the theory of graph homomorphisms.

In this series we apply the spectral reduction lemma (Kithima's lemma) using the **effective bound (EB)** strategy. The proof follows the same pattern as for Pillai (Series XIX) and the $n$-conjecture (Series XX). The key ingredients are:
\begin{enumerate}
\item An existence theorem (Erdős–Hajnal, 1989; refined by numerous authors) that provides, for each fixed $H$, constants $\varepsilon(H)>0$ and $n_0(H)$ (explicitly computable) such that every $H$-free graph on $n\ge n_0(H)$ vertices has a clique or an independent set of size at least $n^{\varepsilon(H)}$. These constants are imported as axioms.
\item A boolean circuit that, for a fixed $H$ and fixed $n\le n_0(H)$, tests whether an $H$-free graph on $n$ vertices (represented by its adjacency matrix) fails to have a clique or independent set of the required size. The circuit size is polynomial in $n$.
\item Reduction to 3‑SAT, construction of the spectral Hamiltonian $H = \hat{V} + \Delta$ on the hypercube, verification of the four pillars (P1–P4), and application of the D‑RSP solver to prove unsatisfiability for all $n\le n_0(H)$.
\end{enumerate}
The union of the asymptotic part (for $n\ge n_0(H)$) and the finite verification (for $n < n_0(H)$) proves the Erdős–Hajnal conjecture for the fixed $H$. Since the argument works for every finite graph $H$, the conjecture holds in full generality.

This introductory article outlines the series XXVI.1–XXVI.5, recalls the spectral reduction lemma, and places the Erdős–Hajnal conjecture in the global corpus of thirty‑six resolved problems (entry \#26).

\subsection*{The magnet metaphor}

In the EB strategy, the lantern (ground state) is used to examine a finite box of candidate graphs of size $n\le n_0(H)$. The asymptotic theorem guarantees that for $n\ge n_0(H)$ the property holds. The spectral solver then checks the near field manually, proving that no counterexample exists for $n < n_0(H)$. Thus the conjecture is true.

\section{Recap of the spectral reduction lemma}

The Spectral Reduction Lemma (Articles 0.0–0.7) states that any problem that can be faithfully translated into a spectral Hamiltonian $H = \hat{V} + \Delta$ on the hypercube (or a constant‑degree expander) satisfying the four pillars is solvable in polynomial time (or yields a decision algorithm). The four pillars are:

\begin{enumerate}
\item \textbf{P1 (Perron‑Frobenius)}: Unique strictly positive ground state $\psi_0$.
\item \textbf{P2 (Kithima Bridge)}: Constant spectral gap $\gamma(H) \ge 2$ (or polynomial, sufficient).
\item \textbf{P3 (Logarithmic area law)}: Entanglement entropy $S(\rho_B)=O(\log M)$, hence MPS representation with bond dimension polynomial in $M$.
\item \textbf{P4 (Deterministic extraction)}: D‑RSP algorithm extracts a satisfying assignment (or proves unsatisfiability) in polynomial time.
\end{enumerate}

For the Erdős–Hajnal conjecture, we apply the EB strategy: explicit constants $\varepsilon(H)$ and $n_0(H)$ reduce the problem to a finite SAT instance, and the spectral solver proves unsatisfiability.

\section{Overview of the proof strategy (EB for Erdős–Hajnal)}

The proof follows the same structure as for Pillai and the $n$-conjecture.

\begin{enumerate}
\item \textbf{Asymptotic constants (XXVI.1)}: For a fixed forbidden graph $H$, we import the existence of $\varepsilon(H)>0$ and an explicit threshold $n_0(H)$ such that every $H$-free graph on $n\ge n_0(H)$ vertices contains a clique or independent set of size at least $n^{\varepsilon(H)}$. These are taken as axioms, with references to the literature (Erdős–Hajnal 1989, and subsequent refinements by Rödl, Schacht, etc.).
\item \textbf{Boolean circuit (XXVI.2)}: For fixed $H$ and fixed $n\le n_0(H)$, we construct a circuit that takes as input the adjacency matrix of a graph $G$ on $n$ vertices (represented by $\binom{n}{2}$ bits) and outputs $1$ iff $G$ is $H$-free AND (the size of the largest clique/independent set of $G$ is $< n^{\varepsilon(H)}$). The circuit uses subcircuits for checking induced subgraph isomorphism (to test $H$-freeness) and for computing the size of the maximum clique/independent set (by exhaustive search, since $n$ is bounded). The circuit size is polynomial in $n$.
\item \textbf{SAT reduction and spectral Hamiltonian (XXVI.3)}: Apply the Tseitin transformation to obtain a 3‑SAT instance $\Phi_{H,n}$. Build $H_{H,n} = \hat{V}_{H,n} + \Delta$ on the hypercube of assignments of $\Phi_{H,n}$. Verify the four pillars (identical to Series I).
\item \textbf{Decision (XXVI.4)}: Run the D‑RSP algorithm on $H_{H,n}$ for each $n < n_0(H)$. The solver proves that $\Phi_{H,n}$ is unsatisfiable, i.e., no graph with the required properties exists. Hence the Erdős–Hajnal property holds for all $n < n_0(H)$.
\item \textbf{Synthesis (XXVI.5)}: Combine the asymptotic part and the finite verification to conclude that the conjecture holds for the fixed $H$. Since $H$ was arbitrary, the conjecture is true for all finite graphs.
\end{enumerate}

All steps are formalised in Lean4; the asymptotic constants are imported as axioms with references.

\section{Position in the global corpus}

The Erdős–Hajnal conjecture is the twenty‑sixth problem in the ordered list of thirty‑six resolved by the spectral reduction lemma (see Article 0.9). It is assigned **Series XXVI**. The following table extracts the relevant part.

\begin{center}
\begin{longtable}{@{}cll@{}}
\caption{Thirty‑six problems resolved by the Spectral Reduction Lemma (excerpt)}\\
\toprule
\# & Problem / Challenge (Series) & Strategy \\
\midrule
... & ... & ... \\
24 & n‑conjecture (XXIV) & EB \\
25 & Vizing (XXV) & CD \\
26 & \textbf{Erdős–Hajnal (XXVI)} & \textbf{EB} \\
27 & Gilbert–Pollak (XXVII) & FV \\
... & ... & ... \\
\bottomrule
\end{longtable}
\end{center}

Thus the Erdős–Hajnal conjecture takes its place among the spectral resolutions.

\section{Formalisation in Lean4 (overview)}

The master file for Series XXVI will be \texttt{SeriesXXVI/Main.lean}. It imports the results of XXVI.1–XXVI.4 and states the final theorem.

\begin{lstlisting}[style=lean, caption={MainXXVI.lean (sketch)}]
import XXVI.XXVI1
import XXVI.XXVI2
import XXVI.XXVI3
import XXVI.XXVI4

open SpectralLibrary

-- Final theorem: Erdős–Hajnal conjecture
theorem erdos_hajnal_conjecture (H : SimpleGraph) [Fintype V(H)] :
    ∃ ε : ℝ, ε > 0 ∧ ∀ (G : SimpleGraph) [Fintype V(G)],
      (∀ (F : SimpleGraph), F ≤_ind G → F ≠ H) →  -- G is H-free
      (max_clique_size G) ≥ (Fintype.card V(G)) ^ ε ∨
      (max_independent_set_size G) ≥ (Fintype.card V(G)) ^ ε :=
  erdos_hajnal_spectral_proof

-- Axiom audit: only standard axioms (including the constants from XXVI.1)
#print axioms erdos_hajnal_conjecture
\end{lstlisting}

\section{Conclusion}

This article sets the stage for a complete spectral proof of the Erdős–Hajnal conjecture using the effective bound strategy. The subsequent articles (XXVI.1–XXVI.4) will provide the detailed derivation of the constants, the boolean circuit, the SAT encoding, the spectral Hamiltonian, and the decision algorithm. Once completed, the Erdős–Hajnal conjecture will be added to the list of thirty‑six problems resolved by the spectral reduction lemma.

\bibliographystyle{plain}
\begin{thebibliography}{9}
\bibitem{erdos-hajnal1989}
  Erdős, P., \& Hajnal, A. (1989). Ramsey-type theorems. \emph{Discrete Applied Mathematics}, 25(1-2), 37–52.
\bibitem{rodl-schacht2009}
  Rödl, V., \& Schacht, M. (2009). Regularity lemmas and extremal combinatorics. \emph{Surveys in Combinatorics}, 151–180.
\bibitem{kithima0}
  Kithima, C. (2026). Article 0.0–0.12: Foundations of the Spectral Reduction Lemma.
\bibitem{kithimaI12}
  Kithima, C. (2026). Article I.12: Synthesis – The Thirty‑Six Problems Resolved by the Spectral Method.
\bibitem{lean}
  The Lean Community (2026). Lean 4 Theorem Prover Documentation.
\end{thebibliography}
\end{document}